% Generated by scripts/report/benchmark_report.py -- do not edit by hand.
\documentclass[10pt,a4paper]{article}
\usepackage[margin=1.9cm]{geometry}
\usepackage{booktabs}
\usepackage{microtype}
\usepackage[T1]{fontenc}
\usepackage{lmodern}
\usepackage[font=small,skip=4pt]{caption}
\usepackage[hidelinks]{hyperref}
\setlength{\parskip}{1pt}
\pagestyle{empty}

\title{\vspace{-1.6cm}\large Model selection for per-receptor odorant activity prediction:\\
cross-validation versus nested cross-validation}
\author{\normalsize Kamran Heydarov}
\date{\normalsize Submitted to Dr.~Emil Sayilov \quad\textbullet\quad 2026-10-10}

\begin{document}
\maketitle
\vspace{-1.25cm}

\paragraph{Question.} The ENodorIM simulator predicts, for an arbitrary molecule, which of the human odorant
receptors it activates. The predictor deployed on the site is one fixed configuration---an $L_2$-regularised
logistic regression per receptor on Morgan fingerprint bits---reported with ordinary $k$-fold cross-validation.
Two questions were conflated by that figure: \emph{how good is this model}, and \emph{how good is the procedure
that chooses a model}. The second is the one that matters before tuning, and it requires the hyperparameter
search to be nested inside the validation. This report answers both on identical folds.

\paragraph{Data.} Human single-compound dose--response and screening experiments from M2OR, reduced to one
consensus label per receptor--molecule pair, giving 654 molecules described by 2048-bit Morgan
fingerprints (radius 2, chirality retained, because measured enantiomers activate receptors differently).
31 receptors have at least eight molecules on each side and are modelled; each has
70--319 tested molecules and 8--134 agonists, so every per-receptor data set is
small and imbalanced.

\paragraph{Protocol.} 5-fold stratified outer folds, repeated 3$\times$ with
different shuffles; out-of-fold scores are pooled across a split before scoring, since a single fold can hold two
positives. \emph{Flat}: hyperparameters chosen on the outer folds and the winner's own score reported---what a
tuning loop reports about itself. \emph{Nested}: hyperparameters chosen by 3-fold inner
cross-validation within each outer training fold, so selection is never graded on data it saw. Fingerprint-bit
filtering is refitted per training fold. The difference between the two columns is the selection bias.

\begin{table}[h]
\centering\small
\setlength{\tabcolsep}{4.5pt}
\begin{tabular}{lrrrrrrrr}
\toprule
Model family & Grid & Flat & Nested & Optimism & Nested AP & $\Delta$ vs dep. & W/L & $p_{\mathrm{Holm}}$ \\
\midrule
Logistic regression, $C=0.1$ (deployed) & 1 & 0.686 & \textbf{0.686} & --- & 0.412 & --- & --- & --- \\
Random forest & 4 & 0.707 & 0.686 & +0.022 & 0.412 & 0.000 & 13/18 & 0.78 \\
Extra trees & 4 & 0.710 & 0.683 & +0.027 & 0.400 & -0.003 & 11/20 & 0.78 \\
Logistic regression, L2 & 6 & 0.707 & 0.677 & +0.030 & 0.405 & -0.009 & 8/23 & 0.26 \\
Bernoulli naive Bayes & 3 & 0.701 & 0.676 & +0.025 & 0.383 & -0.010 & 8/23 & 0.52 \\
Tanimoto $k$-nearest neighbours & 3 & 0.703 & 0.674 & +0.029 & 0.384 & -0.012 & 11/20 & 0.52 \\
Support vector machine, linear & 4 & 0.712 & 0.671 & +0.040 & 0.397 & -0.015 & 9/22 & 0.08 \\
Support vector machine, RBF & 9 & 0.726 & 0.668 & +0.058 & 0.388 & -0.018 & 9/22 & 0.08 \\
Logistic regression, L1 & 5 & 0.709 & 0.657 & +0.051 & 0.373 & -0.029 & 9/22 & 0.20 \\
\bottomrule
\end{tabular}
\caption*{\small Mean per-receptor ROC-AUC over 31 receptors. \emph{Grid}: hyperparameter settings
searched. \emph{Optimism}: flat minus nested. \emph{W/L}: receptors where the family's nested AUC beats or
loses to the deployed model. No family differs from the deployed model after Holm correction over the 8 comparisons (Wilcoxon signed-rank, paired by receptor); 3 reach an uncorrected $p<0.05$ and all are worse, not better.}
\end{table}

\vspace{-0.45cm}
\paragraph{Findings.} (i) No family beats the deployed configuration under nested cross-validation: it is itself the highest-scoring of the 9 at AUC~0.686, matched to three decimals by the random forest and ahead of the remaining 7; no family wins on a majority of receptors.
(ii) Hyperparameter search inflates the flat-CV figure by +0.035 AUC on average---larger than the entire
spread between the 9 families---so a tuned model reported by flat cross-validation would appear to beat the
deployed one while being no better. (iii) Repeating the identical evaluation with a different fold shuffle moves a
single receptor's AUC by $\pm$0.021, again comparable to the between-model differences. The deployed similarity baseline, evaluated by hold-one-out over 654 reference molecules, ranks a true agonist first for 61\% of molecules ($n=138$). 

\paragraph{Conclusion.} On this data set model choice is not the binding constraint: 9 families spanning
linear, kernel, ensemble and instance-based learning fall within the noise of one another, and the untuned
deployed configuration is retained---it carries no selection bias to lose. Per-receptor performance is governed by
the amount of measured data, not the algorithm, which is why the simulator reports each receptor's own
cross-validated AUC and uses a learned score only above AUC~0.70 rather than quoting a single headline number.
The next gain must come from more measured pairs, or from a model that shares one representation across receptors
(a graph encoder with a receptor-sequence embedding), rather than from a wider hyperparameter search.

\paragraph{Reproducibility.} \texttt{python -m scripts.pipeline.benchmark\_models} ($\approx$24~min)
writes \texttt{model\_benchmark.json}, from which this report and the site's benchmark page are generated.
Data: M2OR, \textsc{doi:}10.1093/nar/gkad886. Results: \url{https://kammmran.github.io/ENodorIM/\#/benchmark}.

\end{document}
